\section{Requirements, data and source preparation}
\label{ch:data}

\subsection{From student questions to an evaluation task}

The institution publishes information in several forms: short editorial FAQ answers, regulations whose general rules have exceptions in other articles, timetables that change during the semester, and individual registration states held by authenticated services. Treating all four as interchangeable text would hide differences in authority, freshness and the action an assistant can legitimately perform.

In this thesis, effectiveness concerns the information and guidance returned to the user. A response should address the question, include important conditions, and distinguish an explanation from an action. When the evidence does not support a direct answer, a clarification or referral may be the appropriate outcome. This motivates separate measurements of supported answering and useful routing.

The requirements arise from the observed application and the available question sources. Historical correspondence provides longer descriptions of institutional problems; selectively recorded chatbot questions show how people address a conversational interface. Neither is a complete specification of user needs, and the two are kept apart throughout sampling and analysis.

\begin{table}[htbp]
\centering
\caption{Effectiveness requirements and the evidence available in this thesis. These are analytical requirements, not separately measured acceptance rates.}
\label{tab:requirements}
\begin{tabularx}{\linewidth}{@{}p{0.24\linewidth}XX@{}}
\toprule
Requirement & Observable evidence & Important boundary \\
\midrule
Answer the actual question & Required facts, conditions and response text & References are AI-produced and may be incomplete. \\
Avoid unsupported claims & Source excerpts and material-claim judgments & Quotation occurrence does not establish entailment. \\
Give a useful next step & Clarification, verification or referral in the answer & Guidance is not completed institutional action. \\
Keep evidence traceable & Source identifiers, captured context and prompt & A traceable source can still be outdated. \\
Deliver the response & Recorded completion, disconnection and timing & Completion is not an answer-quality label. \\
Support different interfaces & Shared API and channel adapters & Full web and Discord behavior is outside the replay experiment. \\
\bottomrule
\end{tabularx}
\end{table}

\subsection{The institutional knowledge corpus}

The audited corpus contains 74 FAQ records and 19 documents, split into 1,403 document chunks; together they form the 1,477 indexed vectors of the frozen retrieval corpus. These counts describe the inspected material, not complete coverage of faculty rules or the contents of a later deployment.

FAQ records and document chunks have different roles. An FAQ normally connects a recognizable question to a compact, editorially selected answer, which suits recurring requests with stable wording. A document passage often expresses a rule without repeating the wording of the student's question, and several passages may be needed to establish the rule, its conditions and the procedure. One numerical rejection threshold for both source types may therefore affect them differently. This motivates the threshold experiment, but lowering a document threshold also admits passages that the reranker finds weakly related, so the effect must be examined at the answer level.

The documents repository retains original files and converted Markdown, produced by extracting embedded text or transcribing scanned pages, with a documented requirement to inspect conversions. Conversion quality is treated as part of the evidence chain: an error introduced before indexing can remain perfectly reproducible in retrieval and generation.

Source metadata should not be confused with the text itself. A stable document name supports updates, a title identifies a source, a URL locates the original, and a content hash shows whether submitted text changed. None of these fields establishes institutional applicability. An upload date may differ from the date a rule took effect, and an original PDF and its Markdown rendering have different hashes while describing the same source.

\subsection{Authority, currentness and individual applicability}

Institutional answers require more than finding a sentence that resembles the question. Faculty and university rules can overlap. Transfer provisions appear both in the faculty rules from 2013 and in the university regulations from 2019. They set different minimum shares of credits to be earned at the receiving unit: at least half of the diploma credits in the faculty rules and at least 60\% of the programme credits in the university rules. A 2022 supplement allows the competent body to waive the 60\% minimum when a state of emergency, crisis or war has been declared in the country from which the student transfers~\cite{finki2013study,ukim2019study,ukim2022supplement}. A request containing the word ``transfer'' can therefore match several provisions at different levels of authority. The relevant situation and the rule in force must be identified.

The cited fee schedule applies from the 2024/25 academic year until replacement~\cite{finki2024fees}. The thesis records that source qualification without issuing current fee advice. A passage about a general exemption does not reveal a student's recorded status, so an effective response should explain the public information and identify what remains to be verified.

This affects reference construction. A reference can contain an exact quotation yet overstate its applicability, while a cautious answer may be useful even when the evidence cannot establish an individual outcome. The assessment therefore records required actions, currentness and uncertainty alongside facts.

\subsection{Historical university correspondence}

The supplied university contact archive contains 1,693 incoming-question/reply rows. It was provided as a spreadsheet and examined locally with permission. A reply is evidence of a historical exchange, not a correct reference answer for a new experiment, since its interpretation can depend on the date, the student's circumstances, omitted information, and later changes in rules or systems.

The archive was sampled for several purposes. A 240-case purposive pool emphasized challenging scripts and routes within topics. A separate probability sample of 320 groups served the component study (Section~\ref{sec:component-study}), alongside development and pilot materials. These collections answer different questions and are not merged. The selection differences are substantial: Macedonian Cyrillic appears in 71.94\% of the archive but 39.17\% of the 240-case pool, and romanized Macedonian rose from 7.15\% to 37.08\%. The challenge pool therefore cannot stand in for typical demand, and weighting by broad topic cannot repair selection that favors particular scripts or difficult examples within a topic.

For the controlled Luna experiment, 96 email cases were randomly selected from the 320 email groups of the component study (Section~\ref{sec:component-study}). A separate 20-case email pilot supported execution checks. Selection did not use new answer-quality labels, but these cases had been examined before, so the main email subset is not an untouched holdout.

The 344 component-study email cases (320 test and 24 development) have a median question length of 306 characters. Longer questions can include several conditions, destinations and requested actions, so an answer can cite relevant evidence and still be incomplete. The assessment-control example in Chapter~\ref{ch:methodology} shows why one correct fee statement is insufficient for a question that also asks what documents to submit and where.

\subsection{Questions recorded by the chatbot}

The archived telemetry window contains 699 chatbot requests, but only 172 have recorded message text. Two of those contain no word tokens. Normalizing the remaining 170 turns gives 148 unique question groups, from which the main experiment samples 96 cases. The messages with recorded text are from July, while the broader telemetry window extends beyond that month.

The production label therefore means questions observed in the deployed application, not a representative sample of deployed use. Selective logging may omit periods, paths or request types. Normalization removes exact repeats but does not establish semantic independence: two differently worded questions can concern the same issue, and a repeated short message can mean different things in different conversations.

Only the latest user turn is retained; conversation history is excluded. A follow-up that refers to an earlier answer can therefore be under-specified in replay, and a request for clarification or an assessor disagreement may reflect missing context rather than a retrieval failure.

The production turns have a median length of 31.5 characters, much shorter than the email cases, and include conversational messages as well as institutional questions. Differences between email and production outcomes can reflect length, task type, prior context and selection as well as language or script, so the two frames are reported separately.

\begin{table}[htbp]
\centering
\caption{Question sources and their role in the controlled experiment. Counts refer to distinct stages and must not be summed as independent samples.}
\label{tab:data-flow}
\begin{tabularx}{\linewidth}{@{}p{0.34\linewidth}rX@{}}
\toprule
Collection or stage & Count & Role or limitation \\
\midrule
Contact archive rows & 1,693 & Historical demand; replies are not current ground truth. \\
Component-study email groups & 320 & Sampling frame for the main email subset. \\
Main email questions & 96 & Selected before new answer generation; previously examined. \\
Email execution pilot & 20 & Separate from main estimates. \\
Recorded chatbot requests & 699 & Broader descriptive telemetry window. \\
Requests with recorded text & 172 & Selectively logged July messages. \\
Token-bearing user turns & 170 & Two tokenless records excluded from grouping. \\
Normalized production groups & 148 & Frame for sampling; semantic duplicates can remain. \\
Main production questions & 96 & Latest turns only; history unavailable in replay. \\
\bottomrule
\end{tabularx}
\end{table}

\subsection{Application limits and retained failures}

The chat schema rejects questions longer than 2,000 characters. Four main email questions exceed this limit. The experiment retains each of them in all four policy denominators and records the resulting sixteen input-rejection outcomes, without shortening the questions or generating substitute answers. Each frame and policy therefore contains 96 selected cases, while 92 email and 96 production cases can proceed to generation. Reporting only eligible inputs would answer a narrower question; the full denominator captures this interface boundary, and separate execution counts show how many model calls were possible.

The same principle applies to later failures. A missing, empty or contradictory judgment is not a successful answer and is not dropped from the sample. The methodology preserves both raw records and the rules that produce the final reported category.

\subsection{Evidence separation and responsible reuse}

The knowledge corpus, demand questions and assessment references are separate artifacts. Correspondence and telemetry defined evaluation tasks but were not ingested into the chatbot's knowledge base. References were generated from a bounded source pool, not copied from historical replies. This separates what the assistant knew from what the evaluator expected.

Raw questions, generated answers and case-level review records remain in restricted local research files. The thesis reports aggregate counts and one paraphrased observed example, whose personal wording, category and program are not reproduced. No private dataset is released. The correspondence and telemetry were analyzed with permission, which does not constitute institutional approval of the generated reference facts; the AI-produced references and judgments have not been validated by human assessors.
